\documentclass[10pt]{article}
\usepackage[T1]{fontenc}
\usepackage[utf8]{inputenc}
\usepackage[french]{babel}
\usepackage{amsmath,amssymb,mathrsfs,mathpazo}
\usepackage{geometry}
\geometry{a4paper,margin=2cm}
\usepackage{graphicx,array,tabularx,tabularray,colortbl,multirow,multicol,booktabs,enumitem,pifont,ulem,fancybox,diagbox,cancel}
\usepackage{pstricks,pstricks-add,pst-plot,pst-tree,pst-node,pst-text,pst-func,pst-math,pst-bspline,pst-3dplot,pst-eucl,pst-solides3d}
\let\pscalculate\undefined
\usepackage{pst-bezier}
\usepackage{tikz,pgfplots,tkz-tab}
\usetikzlibrary{arrows,arrows.meta,patterns,calc,positioning,trees,decorations.markings,shapes.geometric,shapes.misc,shapes.arrows,intersections,decorations.pathmorphing,decorations.pathreplacing,backgrounds,shadows,babel}
\pgfplotsset{compat=1.18}
\usepgfplotslibrary{fillbetween}
\usepackage[np]{numprint}
\usepackage{listings,makecell,pas-tableur,twoopt}
\let\barre\undefined
\usepackage[pstricks]{tablvar}
\definecolor{darkgreen}{RGB}{0,100,0}
\providecommand{\red}{\color{red}}
\providecommand{\blue}{\color{blue}}
\providecommand{\green}{\color{green}}
\providecommand{\black}{\color{black}}
\providecommand{\white}{\color{white}}
\providecommand{\gray}{\color{gray}}
\providecommand{\cyan}{\color{cyan}}
\providecommand{\magenta}{\color{magenta}}
\providecommand{\yellow}{\color{yellow}}
\providecommand{\violet}{\color{violet}}
\providecommand{\orange}{\color{orange}}
\providecommand{\euro}{\texteuro}
\providecommand{\EUR}{\texteuro}
\providecommand{\ieme}{\textsuperscript{e}}
\providecommand{\degres}{\ensuremath{^\circ}}
\providecommand{\degree}{\ensuremath{^\circ}}
\providecommand{\vv}[1]{\overrightarrow{#1}}
\providecommand{\Coord}[1]{#1}
\providecommand{\decosix}{}
\providecommand{\starredbullet}{\ensuremath{\bullet}}
\pagestyle{empty}
\setlength{\parindent}{0pt}
\usepackage{adjustbox}\DeclareUnicodeCharacter{2236}{:}\DeclareUnicodeCharacter{2190}{\ensuremath{\leftarrow}}
\providecommand{\R}{\mathbb{R}}
\providecommand{\N}{\mathbb{N}}
\providecommand{\Z}{\mathbb{Z}}
\providecommand{\Q}{\mathbb{Q}}
\providecommand{\D}{\mathbb{D}}
\providecommand{\C}{\mathbb{C}}
\providecommand{\e}{{\mathrm{e}}}
\providecommand{\vect}[1]{\overrightarrow{#1}}
\providecommand{\vectt}[1]{\overrightarrow{\mathrm{#1}}}
\providecommand{\vvt}[1]{\overrightarrow{\mathrm{#1}}}
\providecommand{\np}[2][]{{#2~#1}}
\providecommand{\vv}[1]{{\overrightarrow{#1}}}
\renewcommand{\d}{{\mathrm{d}}}
\providecommand{\degres}{{^\circ}}
\providecommand{\degre}{{^\circ}}
\providecommand{\blue}{\color{blue}}
\providecommand{\red}{\color{red}}
\providecommand{\enskip}{\quad}
\providecommand{\llbracket}{\lbrack\!\lbrack}
\providecommand{\rrbracket}{\rbrack\!\rbrack}
\providecommand{\euro}{{\text{€}}}
\providecommand{\no}{{\text{n°}}}
\providecommand{\interval}[2]{{[#1;#2]}}
\providecommand{\intervFO}[2]{{[#1;#2[}}
\providecommand{\dotfill}{\ldots}

\providecommand{\ds}{\displaystyle}
\providecommand{\textding}[1]{#1}
\providecommand{\Oij}{\(\mathrm{O};\vec{\imath},\vec{\jmath}\)}
\providecommand{\Oijk}{\(\mathrm{O};\vec{\imath},\vec{\jmath},\vec{k}\)}
\providecommand{\Ouv}{\(\mathrm{O};\vec{u},\vec{v}\)}
\providecommand{\intervFF}[2]{\ensuremath{[#1;#2]}}
\renewcommand{\euro}{\ifmmode\text{\texteuro}\else\texteuro\fi}
\DeclareUnicodeCharacter{00B0}{\ifmmode^{\circ}\else\textdegree\fi}
\DeclareUnicodeCharacter{20AC}{\ifmmode\text{\texteuro}\else\texteuro\fi}
% Énoncés : Éducation nationale. Transcription et corrigés : APMEP, Denis Vergès et contributeurs. Découpage et classement : CoopMaths, Sébastien Lozano, Éric Elter et contributeurs. Adaptation de la présentation : Maths Explorer.
% CC BY-SA : https://coopmaths.fr/www/about/
% https://coopmaths.fr/alea/static/e3c/2020/tex/e3c_2020_00_sujet11_4_cor.tex
\begin{document}
\noindent\textit{APMEP / CoopMaths. CC BY-SA. Présentation adaptée par Maths Explorer.}\par\medskip
\subsection*{1. Déterminer le coût moyen quotidien pour la production de 5 m$^3$ d'engrais.}
			
			On a :
			\[
			f(5) = \dfrac{5^2 - 15 \times 5 + 400}{5} = \dfrac{25 - 75 + 400}{5} = \dfrac{350}{5} = 70
			\]
			
			Le coût moyen de 5 m$^3$ d'engrais est 7 000 €.
			
			\subsection*{2. Quels volumes d'engrais faut-il fabriquer pour avoir un coût moyen de production égal à 4 300 € (43 centaines d'euros) ?}
			
			Il faut résoudre l'équation :
			\[
			f(x) = 43 \iff \dfrac{x^2 - 15x + 400}{x} = 43 \iff x^2 - 15x + 400 = 43x \quad (avec \ x \neq 0)
			\]
			\[
			\iff x^2 - 58x + 400 = 0
			\]
			
			Avec $\Delta = 58^2 - 4 \times 400 = 3364 - 1600 = 1764 = 42^2$, les racines sont :
			\[
			x_1 = \dfrac{58 + 42}{2} = 50 \quad \text{et} \quad x_2 = \dfrac{58 - 42}{2} = 8
			\]
			
			En fabriquant 8 ou 50 m$^3$ d'engrais le coût moyen de production est 4 300 €.
			
			\subsection*{3. Pour quel volume d'engrais fabriqué le coût moyen de production est-il minimal ? Déterminer ce coût moyen minimal.}
			
			La fonction $f$ est dérivable sur $[5 ; 60]$ et sur cet intervalle :
			\[
			f'(x) = \dfrac{x(2x - 15) - (x^2 - 15x + 400)}{x^2} = \dfrac{2x^2 - 15x - x^2 + 15x - 400}{x^2} = \dfrac{x^2 - 400}{x^2} = \dfrac{(x + 20)(x - 20)}{x^2}
			\]
			
			$f'(x)$ est du signe de $(x + 20)(x - 20)$ (car $x^2 > 0$) et l'on sait que ce trinôme est positif sauf sur $] - 20 ; 20[$.
			
			Donc sur $[5 ; 20[$, $f'(x) < 0$ et $f$ est donc décroissante et sur $[20 ; 60]$, la fonction $f$ est croissante. \\
			D'où le tableau de variations (avec $f(20) = 25$) :
			\begin{center}
			\begin{tikzpicture}[double distance=2pt]
				\tkzTabInit{$x$/1,$f'(x)$/1,$f(x)$/2}{$5$,$20$, $60$}
				\tkzTabLine{,-,z, +}
				\tkzTabVar{+/$70$,-/$25$, +/$\alpha$}
			\end{tikzpicture}
		\end{center} 
			avec $\alpha=f(60)\approx51,7$\\
			D'après cette étude le minimum est atteint pour $x = 20$ et ce coût moyen minimal est égal à $f(20) = 25$, soit 2 500 €.

\end{document}
